\chapter{KJI--SAXPY 算法与块分解}

\section{KJI 循环次序}

设 $A\in\mathbb{R}^{n\times n}$ 进行原位 LU 分解，暂不写置换。第 $k$ 步先计算乘子

\[
l_{ik}=\frac{a_{ik}}{a_{kk}},
\qquad i=k+1,\ldots,n,
\]

再更新尾矩阵

\[
a_{ij}\leftarrow a_{ij}-l_{ik}a_{kj},
\qquad i,j=k+1,\ldots,n.
\]

KJI 次序把 $k$ 放在最外层、$j$ 放在中层、$i$ 放在最内层：

\begin{lstlisting}
for k = 1:n-1
    A(k+1:n,k) = A(k+1:n,k) / A(k,k)
    for j = k+1:n
        A(k+1:n,j) = A(k+1:n,j)
                     - A(k,j) * A(k+1:n,k)
    end
end
\end{lstlisting}

固定 $k,j$ 后，内层更新为

\[
A(k+1:n,j)\leftarrow
A(k+1:n,j)-A(k,j)A(k+1:n,k),
\]

即向量形式

\[
y\leftarrow y+\alpha x,
\]

对应 BLAS-1 的 SAXPY；双精度实数对应 DAXPY。列主序存储下，内层 $i$ 遍历同一列的连续元素。

对每个 $k$，全部后续列的更新还可以写成秩一形式

\[
A_{k+1:n,k+1:n}\leftarrow
A_{k+1:n,k+1:n}
-A_{k+1:n,k}A_{k,k+1:n}.
\]

该形式对应 BLAS-2 的 GER。SAXPY 和 GER 表示的是同一消元更新的两种接口组织。将多个秩一更新合并为块更新，是稠密分块 LU 和超节点稀疏分解使用 BLAS-3 的基础\cite{davis2006,demmel1999}。

\section{多个秩一更新的合并}

若连续处理 $b$ 个主元 $k=s,\ldots,s+b-1$，这些主元对后续子矩阵的总更新为

\[
\sum_{q=s}^{s+b-1}l_qu_q^T.
\]

把列向量 $l_q$ 和行向量 $u_q^T$ 分别组合成

\[
L_P=
\begin{bmatrix}l_s&l_{s+1}&\cdots&l_{s+b-1}\end{bmatrix},
\]

\[
U_P=
\begin{bmatrix}
u_s^T\\u_{s+1}^T\\\vdots\\u_{s+b-1}^T
\end{bmatrix},
\]

则

\[
\sum_{q=s}^{s+b-1}l_qu_q^T=L_PU_P.
\]

因此 $b$ 次秩一更新可以合并成一次矩阵乘法：

\[
A_{\mathrm{trail}}\leftarrow
A_{\mathrm{trail}}-L_PU_P.
\]

块分解并不改变 LU 的代数结果；它把重复的向量更新重排为矩阵--矩阵运算。

\section{方形分块 LU}

把当前活动矩阵按宽度 $b$ 分块：

\[
A=
\begin{bmatrix}
A_{11}&A_{12}\\
A_{21}&A_{22}
\end{bmatrix},
\qquad A_{11}\in\mathbb{R}^{b\times b}.
\]

\begin{figure}[htbp]
\centering
\begin{tikzpicture}[font=\small]
  \draw[thick] (0,0) rectangle (5.2,5.2);
  \draw[thick] (2.05,0)--(2.05,5.2);
  \draw[thick] (0,3.15)--(5.2,3.15);
  \fill[MumpsBlue!16] (0,3.15) rectangle (2.05,5.2);
  \fill[MumpsCyan!12] (2.05,3.15) rectangle (5.2,5.2);
  \fill[MumpsCyan!12] (0,0) rectangle (2.05,3.15);
  \fill[gray!12] (2.05,0) rectangle (5.2,3.15);
  \node at (1.025,4.175) {$A_{11}$};
  \node at (3.625,4.175) {$A_{12}$};
  \node at (1.025,1.575) {$A_{21}$};
  \node at (3.625,1.575) {$A_{22}$};
  \draw[MumpsRed,very thick] (0,0)--(0,5.2)--(2.05,5.2)--(2.05,0)--cycle;
  \node[MumpsRed] at (1.025,-0.38) {宽度 $b$ 的面板};
\end{tikzpicture}
\caption{块 LU 的方形活动矩阵和面板。}
\label{fig:kji-block-square}
\end{figure}

块 LU 恒等式为

\[
\begin{bmatrix}
A_{11}&A_{12}\\A_{21}&A_{22}
\end{bmatrix}
=
\begin{bmatrix}
L_{11}&0\\L_{21}&I
\end{bmatrix}
\begin{bmatrix}
U_{11}&U_{12}\\0&S_{22}
\end{bmatrix}.
\]

比较分块乘积得到：

\begin{align}
A_{11}&=L_{11}U_{11},\\
U_{12}&=L_{11}^{-1}A_{12},\\
L_{21}&=A_{21}U_{11}^{-1},\\
S_{22}&=A_{22}-L_{21}U_{12}.
\end{align}

对应一次面板迭代的计算顺序为：

\begin{enumerate}
  \item 在面板中分解主元块 $A_{11}$；
  \item 用左三角求解计算 $U_{12}$；
  \item 用右三角求解计算 $L_{21}$；
  \item 用矩阵乘法更新 $A_{22}$。
\end{enumerate}

第 2、3 步对应 TRSM，第 4 步对应 GEMM。标量 KJI 中连续的 DAXPY/GER 更新由最后一步的 GEMM 集中完成。

\section{主元交换与面板更新}

采用部分选主元时，第 $k$ 列的主元位置为

\[
p_k=\argmax_{i\ge k}|a_{ik}|.
\]

搜索范围是当前列的全部活动行，不限于 $b\times b$ 的对角方块。因此，算法中的“面板”是高为活动行数、宽为 $b$ 的长方形列块；\cref{fig:kji-block-square} 中红色区域表示这一列块。

行交换必须同步作用于当前活动矩阵和已经形成的 $L$ 行，并记录到行置换中。若算法还允许列交换，则活动矩阵的列、全局列编号和列置换必须同步更新，分解关系写成

\[
PAQ=LU.
\]

面板内部不能直接使用尚未包含前面板列贡献的旧数据。设当前面板起点为 $s$，已经接受 $t$ 个主元，准备处理 $k=s+t$ 时，当前列先物化为

\[
A(k:n,k)\leftarrow
A(k:n,k)-L(k:n,s:k-1)U(s:k-1,k).
\]

面板结束后，再把全部已接受列对面板外尾矩阵的贡献一次写回。

\section{延迟主元下的分块 LU}

非对称延迟 LU 算子采用 64 列面板。面板内每一步仍先检查当前列；只有当前列全零时，才按算子规定的活动子块或活动列搜索规则继续选取主元。

若搜索发生在面板尚未完成时，必须先落实该面板已经接受的主元对搜索区域的更新。设面板已经接受 $t\le b$ 列，则写回

\[
A_{\mathrm{trail}}\leftarrow
A_{\mathrm{trail}}-L_P(:,1:t)U_P(1:t,:).
\]

写回后，被检查的活动区是当前 Schur 补。若仍不存在可用主元，则以当前位置为 \code{info}，剩余区间作为延迟后缀。

因此，固定面板宽度 $b$ 是一次最多收集的主元数；遇到延迟时，实际面板宽度可以小于 $b$。

\section{对称块更新}

对称不定分解中，一个面板收集 $1\times1$ 和 $2\times2$ 主元块。设面板乘子为 $L_P$，块对角主元为 $D_P$，尾矩阵更新为

\[
A_{22}\leftarrow A_{22}-L_PD_PL_P^T.
\]

定义

\[
W_P=L_PD_P,
\]

则

\[
A_{22}\leftarrow A_{22}-L_PW_P^T.
\]

分块 Bunch--Kaufman 算子用

\[
S=A_{\mathrm{stored}}-L_PW_P^T
\]

表示面板中尚未写回的 Schur 补。主元测试需要某一列或某一行时，以矩阵--向量乘法计算相应部分；面板完成后，以对称秩二 $k$ 更新写回下三角尾矩阵。面板边界不允许把一个 $2\times2$ 主元拆成两部分。

\section{块三角求解}

块化同样用于因子求解。把下三角系统分块：

\[
\begin{bmatrix}
L_{11}&0\\L_{21}&L_{22}
\end{bmatrix}
\begin{bmatrix}Y_1\\Y_2\end{bmatrix}
=
\begin{bmatrix}B_1\\B_2\end{bmatrix}.
\]

前代为

\[
L_{11}Y_1=B_1,
\]

\[
B_2' = B_2-L_{21}Y_1,
\]

\[
L_{22}Y_2=B_2'.
\]

上三角回代同理。单个右端项时，中间更新是矩阵--向量乘法；多个右端项组成矩阵 $B$ 时，更新

\[
B_2' = B_2-L_{21}Y_1
\]

成为矩阵--矩阵乘法。因而，块因子格式既用于分解阶段合并 SAXPY 更新，也用于求解阶段合并多个右端项的三角更新。

\section{运算量与数据访问}

无论采用标量 KJI 还是块 LU，稠密 $n\times n$ LU 的主导浮点运算量均约为

\[
\frac{2}{3}n^3.
\]

二者的区别是运算组织。标量形式反复读取乘子列和后续列并执行向量更新；块形式把 $b$ 个主元的更新合并，使 $L_P$ 和 $U_P$ 的元素在一次矩阵乘法中重复使用。面板内部仍保留主元搜索、交换和细粒度更新，面板外的主要 Schur 补更新则由 TRSM 和 GEMM 完成。

这一变换可以概括为

\[
\underbrace{\sum_{q=s}^{s+b-1}l_qu_q^T}_{b\text{ 次秩一更新}}
\quad\longrightarrow\quad
\underbrace{L_PU_P}_{\text{一次块更新}}.
\]
